% --- 1. INFORMAZIONI GENERALI ---
\section{Informazioni Generali}
\begin{itemize}
	\item \textbf{Luogo/Piattaforma:} {Server Discord}
	\item \textbf{Data:} 2026-04-07
	\item \textbf{Orario di inizio:} 14:30
	\item \textbf{Orario di fine:} 16:45
	\item \textbf{Responsabile:} {Sofia De Blasi}
	\item \textbf{Scriba:} {Sofia De Blasi}
\end{itemize}

\vspace{0.5cm}
\textbf{Partecipanti presenti:}
\begin{table}[ht]
	\centering
	\renewcommand{\arraystretch}{1.2}
	\begin{tabularx}{\textwidth}{X l}
		\toprule
		\textbf{Nome e Cognome} & \textbf{Matricola} \\
		\midrule
		Sofia De Blasi & 2111014 \\
		Roberto Mariano Doroftei & 2111031 \\
		Filippo Idiometri & 2035617 \\
		Francesco Marcon & 2110990 \\
		Armando Moda Scarati & 2082864 \\
		Anton Ricardo Rupi & 2054304 \\
		\bottomrule
	\end{tabularx}
\end{table}

\vspace{0.5cm}
\textbf{Partecipanti assenti:}
\begin{table}[ht]
	\centering
	\renewcommand{\arraystretch}{1.2}
	\begin{tabularx}{\textwidth}{X l}
		\toprule
		\textbf{Nome e Cognome} & \textbf{Matricola} \\
		\midrule
		Nessuno & - \\
		\bottomrule
	\end{tabularx}
\end{table}

% --- 2. ORDINE DEL GIORNO ---
\section{Ordine del Giorno}
Gli argomenti previsti per la riunione odierna sono:
\begin{enumerate}
	\item Definizione del workflow per la gestione dei task e del tracciamento ore
	\item Avvio della redazione del Piano di Progetto (PdP)
	\item Pianificazione del primo sprint e assegnazione dei ruoli
	\item Avvio della redazione dell'Analisi dei Requisiti (AdR)
	\item Definizione del sistema centralizzato di misurazione ore e issue tracking
\end{enumerate}

% --- 3. RESOCONTO E DISCUSSIONE ---
\section{Resoconto della Riunione}
\subsection{Gestione dei task e workflow operativo}
Il gruppo ha discusso come gestire attività che non producono modifiche dirette alla repository (es.\ compilazione di form Google, comunicazioni interne). Si è deciso di utilizzare una sezione generica nel form per feature/altro. Per le attività non rendicontabili ai fini del progetto (es.\ studio individuale autonomo, attività di formazione personale) il tempo inserito sarà pari a 0. È stato inoltre stabilito di mantenere aggiornato il GitHub Project con stati progressivi per ogni issue. L'apertura dei branch avverrà manualmente da parte di ciascun membro a partire dalla relativa issue, senza automazioni. Il tempo effettivo impiegato per ogni task verrà inserito direttamente nella cella dedicata all'interno del foglio Excel generato dal form Google.
 
\subsection{Avvio della redazione del Piano di Progetto}
Il gruppo ha definito la struttura iniziale del Piano di Progetto. Le sezioni da redigere in questo sprint sono: analisi dei rischi (suddivisa in tre categorie: tecnologici, individuali/organizzativi, relativi al prodotto), modello di sviluppo (Scrum con rotazione bisettimanale), pianificazione a lungo termine (RTB: marzo–luglio 2026; PB: luglio–agosto 2026) e struttura primaria della pianificazione a breve termine, da popolare in modo incrementale sprint per sprint. Data la disponibilità limitata di un singolo redattore per questo documento nel primo sprint, si procederà con l'analisi dei rischi e la struttura delle restanti sezioni.
 
\subsection{Pianificazione del primo sprint e assegnazione dei ruoli}
Il primo sprint è stato definito con durata dal 07 al 14 aprile 2026. I ruoli assegnati per questo sprint sono: Responsabile (Sofia De Blasi), Amministratore (Francesco Marcon), Verificatore (Filippo Idiometri), Analisti (Roberto Mariano Doroftei, Armando Moda Scarati e Anton Ricardo Rupi). 
Dato che attualmente i ruoli di progettista e programmatore non sono ancora necessari, si è deciso che tre membri siano contribuiscano come analisti, più specificatamente Roberto Mariano Doroftei e Armando Moda Scarati si occuperanno dell'AdR, mentre Anton Ricardo Rupi del PdP. 
Si è stabilito, inoltre, che il Responsabile redige il verbale e l'Amministratore gestisce la creazione delle issue e controlla, prima della fine di ogni sprint, che i tempi effettivi delle task nel foglio Excel siano stati compilati correttamente (il workflow prevede che la compilazione avvenga dopo il merge e la chiusura della issue).
 
\subsection{Avvio della redazione dell'Analisi dei Requisiti}
Il gruppo ha stabilito la struttura iniziale dell'AdR. Gli attori del sistema sono stati identificati come: Utente (attore primario), LLM e, opzionalmente, Database/File System (attori secondari). Il template dei Casi d'Uso adottato seguirà il pattern con i seguenti campi: codice identificativo, titolo, attori, descrizione, precondizione, postcondizione, scenario principale e scenari alternativi/estensioni. Lo strumento scelto per i diagrammi UML è \href{https://lucid.co/lucidspark}{\underline{Lucidspark}}. 

I casi d'uso sono stati organizzati in tre macroaree funzionali:
 
\begin{itemize}
	\item \textbf{Macroarea 1 -- Gestione Editor e Markdown:} creazione di una nuova nota vuota, scrittura testo di base, applicazione delle formattazioni (grassetto, corsivo, sottolineato, barrato), inserimento titoli/intestazioni, inserimento link ipertestuali, inserimento elenchi puntati e numerati, selezione del testo, undo/redo, modifica della modalità di visualizzazione (Solo Editor, Solo Render, Split View).
	\item \textbf{Macroarea 2 -- Interazione AI (LLM):} generazione riassunto, traduzione, riscrittura del testo, generazione automatica da prompt (Distant Writing), critica del testo secondo i sei cappelli di De Bono (Bianco, Rosso, Nero, Giallo, Verde, Blu), inserimento della proposta AI nel documento, rifiuto/annullamento della proposta AI, gestione errori di connessione/timeout LLM. Tra i requisiti opzionali: ChatBot con messaggi precompilati.
	\item \textbf{Macroarea 3 -- Gestione I/O e Persistenza:} download della nota in formato Markdown, upload di una nota locale, esportazione (PDF, HTML o JSON). Tra i requisiti opzionali: salvataggio su Database remoto, caricamento da Database remoto, creazione di collegamenti tra note nel Database.
\end{itemize}
 
Per ogni macroarea, i requisiti opzionali saranno raggruppati in una sezione dedicata. 

Per questo sprint l'obiettivo è completare fino alla sezione 3.2 (introduzione) dell'AdR e il primo caso d'uso (Creazione di una nuova nota vuota), lavorando in coppia per calibrare i tempi e la convenzione da adottare. Sarà predisposto un file condiviso per raccogliere dubbi da sottoporre a Zucchetti per chiarimenti sui requisiti.
 
\subsection{Sistema centralizzato di misurazione ore e issue tracking}
Si è definita una struttura centralizzata basata su un unico Google Form con i seguenti campi aggiuntivi rispetto a quelli precedentemente introdotti per ogni task: ruolo con nome dell'assegnatario (es.\ Verificatore, Amministratore, Responsabile) che seguirà l'andamento della issue, tempo previsto dal Verificatore per verificare il prodotto del redattore (proporzionale a quello del redattore), tempo fisso stimato per Responsabile e Amministratore per seguirne l'operato (proporzionale a quello del redattore), tempi effettivi a consuntivo, numero di sprint di riferimento.

% --- 4. DECISIONI PRESE ---
\section{Decisioni Prese}

\renewcommand{\arraystretch}{1.3}
\vspace{-0.3cm}
\noindent
\begin{longtable}{c p{12cm}}
	\toprule
	\textbf{Codice} & \textbf{Decisione} \\
	\midrule
	\endfirsthead
 
	\toprule
	\textbf{Codice} & \textbf{Decisione} \\
	\midrule
	\endhead
 
	\midrule
	\multicolumn{2}{r}{\textit{Continua nella pagina successiva}} \\
	\midrule
	\endfoot
 
	\bottomrule
	\endlastfoot
 
	\textbf{VI-1.1} & I task non relativi a modifiche di repository vengono gestiti tramite una sezione generica nel Google Form (feature/altro). \\[4pt]
	\textbf{VI-1.2} & Le attività non rendicontabili (es.\ studio individuale autonomo) vengono registrate con tempo pari a 0. \\[4pt]
	\textbf{VI-1.3} & Il GitHub Project viene mantenuto aggiornato con stati progressivi per ogni issue. \\[4pt]
	\textbf{VI-1.4} & L'apertura dei branch avviene manualmente da ciascun membro a partire dalla issue, senza automazioni. \\[4pt]
	\textbf{VI-1.5} & Il tempo effettivo viene inserito nella cella dedicata del foglio Excel dopo il merge e la chiusura della issue. \\[4pt]
	\textbf{VI-1.6} & Il Piano di Progetto verrà redatto con le sezioni: analisi rischi (tre categorie), modello di sviluppo Scrum con rotazione bisettimanale, pianificazione a lungo termine (RTB marzo–luglio, PB luglio–agosto) e struttura della pianificazione a breve termine incrementale. \\[4pt]
	\textbf{VI-1.7} & Il primo sprint ha durata dal 07/04/2026 al 14/04/2026. \\[4pt]
	\textbf{VI-1.8} & I ruoli del primo sprint sono: Responsabile (Sofia De Blasi), Amministratore (Francesco Marcon), Verificatore (Filippo Idiometri), 1°Analista (Roberto Mariano Doroftei), 2° Analista (Armando Moda Scarati), 3° Analista--PdP (Anton Ricardo Rupi). \\[4pt]
	\textbf{VI-1.9} & Il Responsabile redige il verbale interno di ogni riunione. \\[4pt]
	\textbf{VI-1.10} & L'Amministratore gestisce la creazione delle issue e verifica prima della chiusura di ogni sprint che i tempi effettivi nel foglio Excel siano stati compilati. \\[4pt]
	\textbf{VI-1.11} & Gli attori del sistema per l'AdR sono: Utente (primario), LLM e opzionalmente Database/File System (secondari). \\[4pt]
	\textbf{VI-1.12} & Lo strumento per i diagrammi dei Casi d'Uso è Lucidspark. \\[4pt]
	\textbf{VI-1.13} & I Casi d'Uso sono organizzati in tre macroaree (Editor e Markdown, Interazione AI, Gestione I/O e Persistenza); per ogni macroarea i requisiti opzionali sono raggruppati in una sezione dedicata. \\[4pt]
	\textbf{VI-1.14} & L'obiettivo dell'AdR per questo sprint è completare fino alla sezione 3.2 e il primo caso d'uso (Creazione nuova nota vuota), lavorando in coppia per calibrare tempi e convenzioni. \\[4pt]
	\textbf{VI-1.15} & Viene predisposto un file condiviso per raccogliere dubbi da sottoporre a Zucchetti per chiarimenti sui requisiti. \\[4pt]
	\textbf{VI-1.16} & Il sistema di tracciamento ore è centralizzato in un unico Google Form con campi per: tempo previsto, ruolo, tempi proporzionali per Verificatore/Responsabile/Amministratore, tempi effettivi e numero sprint. \\
\end{longtable}

% --- 5. AZIONI (TODO) ---
\section{Azioni da Svolgere (To-Do)}

\renewcommand{\arraystretch}{1.3}
\begin{longtable}{c c p{4.5cm} p{3cm} l}
	\toprule
	\textbf{Codice} & \textbf{Rif. Decisione} & \textbf{Descrizione Task} & \textbf{Assegnatario} & \textbf{Scadenza} \\
	\midrule
	\endfirsthead
 
	\toprule
	\textbf{Codice} & \textbf{Rif. Decisione} & \textbf{Descrizione Task} & \textbf{Assegnatario} & \textbf{Scadenza} \\
	\midrule
	\endhead
 
	\midrule
	\multicolumn{5}{r}{\textit{Continua nella pagina successiva}} \\
	\midrule
	\endfoot
 
	\bottomrule
	\endlastfoot
 
	\href{https://github.com/Synergo-Unit-UniPD/Documentary-Repo/issues/103}{\underline{\#103}} & VI-1.9 & Redazione verbale interno riunione del 07/04/2026 & Sofia De Blasi & 2026-04-14 \\[4pt]
	\href{https://github.com/Synergo-Unit-UniPD/Documentary-Repo/issues/104}{\underline{\#104}} & VI-1.14 & AdR: stesura fino a sezione 3.2 + primo caso d'uso (Creazione nuova nota vuota) & Armando M.\ Scarati, Roberto M.\ Doroftei & 2026-04-14 \\[4pt]
	\href{https://github.com/Synergo-Unit-UniPD/Documentary-Repo/issues/105}{\underline{\#105}} & VI-1.6 & PdP: analisi rischi + struttura primaria delle sezioni & A.\ Ricardo Rupi & 2026-04-14 \\[4pt]
	\href{https://github.com/Synergo-Unit-UniPD/Documentary-Repo/issues/106}{\underline{\#106}} & VI-1.16 & Aggiornamento Google Form con nuovi campi per gestione ruoli (Amministratore, Responsabile, Verificatore) & Francesco Marcon & 2026-04-14 \\[4pt]
	\href{https://github.com/Synergo-Unit-UniPD/Documentary-Repo/issues/107}{\underline{\#107}} & VI-1.12 & Studio individuale della teoria sui Casi d'Uso & Tutti i membri & 2026-04-14 \\
\end{longtable}

% --- 6. PROSSIMA RIUNIONE ---
\section{Prossima Riunione}
\begin{itemize}
	\item \textbf{Data:} 2026-04-14
	\item \textbf{Ora:} 14.30
	\item \textbf{OdG preliminare:} 
	\begin{itemize}
		\item Revisione e verifica dei task del primo sprint;
		\item allineamento sulle convenzioni adottate per i Casi d'Uso dell'AdR;
		\item avanzamento del Piano di Progetto;
		\item avanzamento dell'Analisi dei Requisiti;
		\item aggiornamento Norme di Progetto;
		\item pianificazione del secondo sprint e assegnazione dei ruoli.
	\end{itemize}   
\end{itemize}